\documentclass[11pt]{article}

\usepackage[margin=1in]{geometry}
\usepackage{amsmath, amssymb, amsthm}
\usepackage{booktabs}
\usepackage{array}
\usepackage{listings}
\usepackage{xcolor}
\usepackage{hyperref}

\hypersetup{colorlinks=true, linkcolor=blue!50!black, urlcolor=blue!50!black}

\lstset{
  basicstyle=\ttfamily\small,
  keywordstyle=\color{blue!60!black}\bfseries,
  commentstyle=\color{green!40!black}\itshape,
  stringstyle=\color{orange!70!black},
  showstringspaces=false,
  breaklines=true,
  frame=single,
  rulecolor=\color{black!25},
  language=Python,
}

\newtheorem{theorem}{Theorem}
\newtheorem{lemma}{Lemma}
\newtheorem{proposition}{Proposition}
\theoremstyle{definition}
\newtheorem{definition}{Definition}

\title{Prime Certificate:\\ The Mod-6 Wheel, Base 60, and the Totient}
\author{Shaimaa Soltan}
\date{\today}

\begin{document}
\maketitle

\begin{abstract}
This note records the structural chain connecting a period-6 sawtooth,
the base-60 numeral inheritance, wheel factorization, and Euler's totient.
The central result is that a single function, $\varphi$, plays two roles:
it measures the \emph{width of a prime-sieving wheel} and it encodes the
\emph{primality of an individual integer}.
\end{abstract}

%==================================================================
\section{Background: the mod-6 wheel}
%==================================================================

Every integer occupies one of six lanes modulo $6$. Lanes
$0,2,4$ are even and lane $3$ is divisible by $3$, so beyond the base
primes $2$ and $3$ only lanes $1$ and $5$ can hold primes:
\[
p > 3 \;\Longrightarrow\; p \equiv 1 \text{ or } 5 \pmod 6,
\qquad\text{i.e.}\qquad p = 6m \pm 1 .
\]
These two survivors are exactly the fold (reset) positions of the two
sawtooth waves
\[
B = \bigl((x - \tfrac{6}{8}) \bmod 6\bigr)\cdot 2,
\qquad
D = \bigl((x + 2 - \tfrac{6}{8}) \bmod 6\bigr)\cdot 2 ,
\]
which reset at $x \bmod 6 = 1$ and $x \bmod 6 = 5$ respectively.

%==================================================================
\section{Base-60 inheritance}
%==================================================================

Base $60$ factors as $60 = 6 \times 10$: a mod-$6$ ``tens'' wheel geared to
a base-$10$ ``units'' wheel. We compute in base $10$ but time and angle
keep the base-$60$ rollover, whose reset is precisely the $x \bmod 6 = 5$
event of the wheel above. The sieving power of $60$ lives entirely in its
squarefree kernel $30 = 2\cdot 3\cdot 5$.

%==================================================================
\section*{\textbf{Details: The Totient Relation}}
\addcontentsline{toc}{section}{Details: The Totient Relation}
%==================================================================
\markboth{Details: The Totient Relation}{}

Euler's totient $\varphi$ is the hinge of this entire construction. It
appears in \emph{two} distinct roles: once attached to the wheel modulus
$M$, and once attached to the integer $N$ under test.

\begin{definition}[Totient]
For a positive integer $n$,
\[
\varphi(n) \;=\; \bigl|\{\, k : 1 \le k \le n,\ \gcd(k,n)=1 \,\}\bigr|,
\]
the count of integers up to $n$ that are coprime to $n$.
\end{definition}

\subsection*{Role 1 --- $\varphi(M)$ is the wheel width}

For a wheel with modulus $M$, the spokes are exactly the residues coprime
to $M$, so the number of spokes is $\varphi(M)$ and the surviving density
of the number line is $\varphi(M)/M$. This density is governed by the
multiplicative product formula.

\begin{theorem}[Euler product formula]
\[
\varphi(M) \;=\; M \prod_{p \mid M} \Bigl(1 - \frac{1}{p}\Bigr),
\]
the product ranging over the \emph{distinct} primes dividing $M$.
Consequently the wheel density is
\[
\frac{\varphi(M)}{M} \;=\; \prod_{p \mid M}\Bigl(1 - \frac1p\Bigr).
\]
\end{theorem}

Each factor $(1-1/p)$ is one prime's worth of lanes struck out. This one
equation explains every empirical observation about the wheels:

\begin{center}
\begin{tabular}{r r c c l}
\toprule
$M$ & $\varphi(M)$ & $\varphi(M)/M$ & $\prod(1-1/p)$ & distinct primes of $M$\\
\midrule
$6$    & $2$   & $0.3333$ & $0.3333$ & $\{2,3\}$\\
$30$   & $8$   & $0.2667$ & $0.2667$ & $\{2,3,5\}$\\
$60$   & $16$  & $0.2667$ & $0.2667$ & $\{2,3,5\}$ \quad(no gain over $30$)\\
$210$  & $48$  & $0.2286$ & $0.2286$ & $\{2,3,5,7\}$\\
$2310$ & $480$ & $0.2078$ & $0.2078$ & $\{2,3,5,7,11\}$\\
\bottomrule
\end{tabular}
\end{center}

\noindent Two consequences follow immediately:
\begin{itemize}
  \item \textbf{$60$ gains nothing over $30$:} they share the distinct
        primes $\{2,3,5\}$, so the product---and hence the density---is
        identical. The extra factor of $2$ in $60=2^2\cdot3\cdot5$ never
        enters the product, which counts only \emph{distinct} primes.
  \item \textbf{$210$ does gain:} it introduces a genuinely new factor
        $(1-1/7)$, lowering the density from $0.2667$ to $0.2286$.
\end{itemize}

\subsection*{Role 2 --- $\varphi(N)$ is a primality signature}

For the integer $N$ passed to the scan, the totient encodes primality
directly:
\[
N \text{ is prime} \;\Longleftrightarrow\; \varphi(N) = N-1,
\qquad\qquad
\varphi(p^k) = p^k - p^{k-1}.
\]
A prime shares a factor with nothing below it, so all $N-1$ smaller
integers are coprime---the maximal possible totient. Every composite
satisfies $\varphi(N) < N-1$.

\subsection*{The bridge --- Euler's theorem}

The two roles are joined by Euler's theorem, which also underlies the
fast primality tests (Fermat, Miller--Rabin) that outperform trial
division for large $N$.

\begin{theorem}[Euler]
If $\gcd(a,N)=1$ then
\[
a^{\varphi(N)} \equiv 1 \pmod{N}.
\]
For prime $p$ this specialises, via $\varphi(p)=p-1$, to Fermat's little
theorem $a^{\,p-1}\equiv 1 \pmod p$.
\end{theorem}

Miller--Rabin checks whether $N$ \emph{violates} this congruence: a
violation certifies compositeness with a single modular exponentiation,
whereas the wheel must trial-divide up to $\sqrt N$.

\subsection*{Summary}

\[
\boxed{\;
\underbrace{\varphi(M)}_{\text{sieve width}}
\;=\; M\!\prod_{p\mid M}\!\Bigl(1-\tfrac1p\Bigr)
\qquad\text{and}\qquad
\underbrace{\varphi(N)=N-1}_{\text{primality signature}}
\;\Longleftrightarrow\; N \text{ prime},
\;}
\]
the same function serving once as a wheel's width and once as an integer's
primality signature, bridged by $a^{\varphi(N)}\equiv 1 \pmod N$.

%==================================================================
\section{Reference implementation}
%==================================================================

The mod-$210$ wheel ($210 = 2\cdot3\cdot5\cdot7$, $\varphi(210)=48$ spokes)
used as a trial-division primality scanner:

\begin{lstlisting}
from math import gcd, isqrt

WHEEL210 = [r for r in range(210) if gcd(r, 210) == 1]   # 48 spokes

def wheel_scan_210(N):
    """Smallest factor of N via a mod-210 wheel, or None if N is prime."""
    if N < 2:
        return None
    for p in (2, 3, 5, 7):            # base primes the wheel skips
        if N == p:  return None
        if N % p == 0: return p
    limit = isqrt(N)
    m = 0
    while 210*m + 1 <= limit:
        for r in WHEEL210:
            n = 210*m + r             # actual candidate = block offset + spoke
            if n == 1:                # skip trivial divisor
                continue
            if n > limit:
                break
            if N % n == 0:
                return n              # real factor -> composite
        m += 1
    return None                       # nothing up to sqrt(N) -> prime

def is_prime(N):
    return N >= 2 and wheel_scan_210(N) is None
\end{lstlisting}

\end{document}
